%==============================================================================
% Synthesis essay -- the umbrella analysis of the Phase-2 v3 results.
% Written as prose (essay), not bullets. Compiles with a base TeX install.
%   pdflatex synthesis_essay.tex
% This is the team's month-of-analysis narrative; the four paper skeletons
% (paperA..paperD) are the publishable decompositions of it.
%==============================================================================
\documentclass[11pt]{article}
\usepackage[margin=1in]{geometry}
\usepackage{amsmath}
\usepackage{microtype}
\usepackage{booktabs}
\usepackage{hyperref}

\title{From Five Terabytes to a Single Scalar:\\
What a Coarse Memory-Liveness Signal Can and Cannot Tell Us}
\author{VM Memory-Signal Capture Project --- Research Team}
\date{Synthesis of the Phase-2 v3 results}

\begin{document}
\maketitle

\begin{abstract}
This essay reports what we learned from a 132-cell campaign that observes a
virtual machine's memory purely through a single scalar per snapshot, the
Active Page Fraction (APF). We argue that the central result of the project is
not any one classifier or detector score but a coherent trade: by collapsing a
per-page, two-dimensional change signal into one number, we lose the ability to
identify a workload instance and keep only a \emph{weak} signal of its
\emph{family} (bursty vs.\ continuous) once leakage is controlled. An apparent
perfect family separation, we show, is an artifact of features the pipeline
computes per family; the leakage-free accuracy is barely above baseline. We
trace this trade through four results -- capture cost, analyzer tuning, family
classification, and phase segmentation -- and we are deliberate about the
project's four honest negative findings, each of which we believe is more
instructive than the positive number it sits beside.
\end{abstract}

\section{The premise and the price}

The project asks a narrow question with broad consequences: how much can a
hypervisor learn about a guest by watching only how much of its memory changes
over time? The faithful way to watch memory change is per page. An earlier
iteration of this project (``Pipeline A'') did exactly that, computing a cosine
and a Hamming similarity for every 4\,KiB page between consecutive full-memory
snapshots. For a 1\,GiB guest that is roughly $262{,}144$ pages per snapshot and
a two-dimensional matrix of shape $[\,T \text{ snapshots} \times N \text{ pages}\,]$.
The fidelity is real, and so is the bill: retaining the dumps needed to recompute
that matrix costs on the order of 60\,GiB for a single long cell and a projected
$\approx$5\,TiB for a 90-cell campaign (never fully retained) --- more than the
host could hold.

The Active Page Fraction is the project's answer to that bill. It is one number
per snapshot pair: the fraction of guest pages whose bytes change at all,
\[
  \mathrm{APF} = \frac{1}{N}\sum_{p=1}^{N}\mathbf{1}\big[\exists j:\,a_{p,j}\neq b_{p,j}\big].
\]
A single-bit flip and a full-page rewrite count the same; that coarsening is the
whole point. The streaming capture architecture computes each pair's APF in an
asynchronous helper, appends it to a trajectory, and---in its delete-as-you-go
mode---removes the older dump, so the persistent footprint of a 90-cell campaign
is about 4.7\,MiB, an estimated reduction of order a million-fold (a projection
against the never-retained 2-D campaign). One honest caveat belongs here: the v3
run was captured with \texttt{keep\_dumps} enabled, which disables the delete
step, so dumps were retained until a single end-of-cell sweep and transient disk
scaled with the trajectory rather than holding to the two-to-three-gigabyte
design target; validating that cap needs a delete-as-you-go re-run. The compute
drops in step: a memory-mapped byte difference takes 300--500\,ms against
1.5--3\,s for the per-page Rust kernel.

Capturing the v3 dataset taught us something we did not design for, and it is the
first honest finding. We requested a 500\,ms sampling interval, and the guest did
idle for about that long ($\approx$0.52\,s) between snapshots --- but the host
snapshot cycle dwarfs it and is highly variable: a median of about 9.75\,s per
snapshot, ranging from 2 to 24\,s across the 128 usable cells. The per-duration
snapshot counts are spread to match --- 120-second cells a median of 14 (range
5--43), 300-second cells 26 (12--148), 600-second cells 53 (28--100) --- so the
tidy ``7/20/28'' we first quoted are near the per-duration minima, not typical
values. The cost is dominated by \texttt{pmemsave} itself, a median of 7.06\,s
and about 81\% of the cycle, with suspend ($\approx$0.35\,s) and resume
($\approx$0.05\,s) negligible and the asynchronous read-back off the critical
path; the VM-pause fraction is correspondingly high but not uniform, a median of
0.946 (as low as 0.74). The practical consequence reframes the entire ``how often
should we sample'' question: because the snapshot cycle dwarfs the 500\,ms knob,
tuning the interval below it is inert. The right question is not how often to
sample but how cheaply one can snapshot --- that is, how to shrink
\texttt{pmemsave}.

\section{Tuning as certification, not maximization}

With a signal in hand, the analyzer needs a window and a hop. The tempting move
is to sweep both and report the configuration with the best score. We resisted
it. Instead we built an acceptance-gated search: each candidate window/hop is
judged by five named gates --- stationarity, spectral coverage, a defining
metric (cepstral signal-to-noise for the bursty workloads, coefficient of
variation for the steady ones), a hop-validity condition, and a window-count
floor --- plus a guard against regressing a prior baseline.

The gates earn their keep precisely because they disagree with one another. Ten
of the eleven workloads select a window/hop of $(8,4)$; only the bimodal
hash-table workload prefers $(32,8)$. Five workloads --- the steady
micro-benchmarks that are not bimodal --- pass every gate cleanly, with
coefficients of variation between 0.06 and 0.26, comfortably under the ceiling.
The hash-table workload passes all five gates but is held back by the regression
guard, a near-miss rather than a metric failure. And then the bursty
``phasic'' family produces the second honest finding. Every one of those five
workloads \emph{passes} the defining-metric gate --- their cepstral SNR ranges
from 4.78 to 6.35\,dB, all above the 4.5\,dB bar (though on some cells that SNR
rests on only one to three analysis windows, where the estimate is fragile) ---
yet every one \emph{fails}
the spectral-coverage gate, with coverage near 0.20 against a threshold of 2.0.
The rhythm is unmistakably there; we simply did not capture enough periods of it
at the achievable cadence. A single maximized score would have hidden this. The
gate framework states it plainly: the signature is present, and it is
under-sampled. That is an actionable sentence --- it tells the next capture to
buy more periods, not a different window.

\section{What survives the coarsening: less than it first appears}

If APF throws away per-page detail, what workload information is left? We posed
two tasks. The first is family classification: is the workload bursty
(``phasic'', the synthetic reversible-XOR sandboxes that emulate the write
rhythm of encryption in passes, plus a metadata scanner) or continuous
(``steady'', the memory-pressure micro-benchmarks)? The second is instance
identification: which of the eleven specific workloads is this? We answered both
with ten APF-derived features and a cross-validation that leaves an entire
replicate out.

The first answer looked spectacular and was, on inspection, an illusion. A
RandomForest separates phasic from steady with 100\% accuracy over 118 cells, a
confusion matrix of $[[54,0],[0,64]]$ without a single error. But a leakage audit
deflates it completely: a \emph{single} feature, \texttt{coverage\_ratio} ---
which the analyzer assigns as a constant per family, 0.133 to steady cells and
0.200 to phasic ones --- reproduces the same 100\% on its own. The model was
rediscovering a labelling rule the pipeline had baked into the features, not a
property of memory behavior. When we remove the three family-conditional
features (\texttt{coverage\_ratio}, the steady-only \texttt{cv\_workingset}, and
the phasic-only \texttt{f1\_phase}), accuracy falls to 0.644 against a 0.542
majority baseline. That residual gap is weak and, to be candid, not established:
with only about 33 independent workload-duration groups (the four replicates are
not independent), it does not reach significance, so the honest reading is ``at
most a weak hint, indistinguishable from baseline at the group level'' rather
than a demonstrated signal. We also retract the first draft's stated mechanism: cepstral SNR
does \emph{not} separate the families (alone it scores 0.508, chance), because
steady SNR spans 2.6--8.0\,dB and overlaps the phasic range --- the steady
hash-table workload in fact has the highest SNR of all.

Instance identification is weaker still. Eleven-way accuracy plateaus at 33.1\%
for the RandomForest --- and the right comparator is not the 9.1\% random rate
but the majority-class and single-feature baselines, against which the margin is
modest. The errors are almost entirely \emph{within} a family; the metadata
scanner is the best-separated class (F1 0.60) while the working-set and mmap
sweeps are nearly indistinguishable (F1 0.10--0.12). The block-diagonal-by-family
structure inherits the same leakage, so the trustworthy part is the
\emph{within}-family confusion, and it says APF cannot tell two continuous
memory-pressure patterns apart. The honest summary: a single scalar carries a
coarse, weak hint of workload family and essentially no instance identity.

\section{Segmentation, and the limit of ground truth}

The last result asks an online question: place phase boundaries where a workload
changes regime, and, just as importantly, place none when it does not. We use a
windowed CUSUM detector over window means, with a median-absolute-deviation
robust scale and a single drift knob, specialized by the family label that the
classifier provides. Steady workloads keep the conservative drift; phasic
workloads use a smaller one to enter the detection band on a shallow rhythm.

Ten of the eleven workloads pass acceptance, which is 120 of 132 cells. For the
steady family the criterion is restraint: a workload passes if it stays
stationary and emits essentially no spurious boundaries, and five of six do. The
exception, the page-fault-density workload, injects spurious boundaries and
fails. For the phasic family we intended to validate against the workloads'
emitted phase markers --- and here we hit the project's sharpest honest finding.
Not one phasic cell is marker-rich in the sense the detector needs: the markers
describe workload-internal events, such as the encryption of an individual file,
that never surface as a mean shift in APF at this cadence, a problem compounded
by a wall-clock skew on the capture host. Marker-aligned F1 is therefore not
merely low but structurally unavailable. We fall back to validating phasic
segmentation by the plausibility of the boundary count --- it must land in a
sensible band --- and on that basis all five phasic workloads pass. We are
careful to label this for what it is: the strong claim, that we segment phases
correctly, is not yet supported. What we have is a robust detector that knows
when to stay silent, and a precise statement of what the next capture must fix
so that ``correct'' can be measured at all.

\section{Why four papers, and what they share}

These four results decompose naturally into four focused papers --- the capture
method and its cost model; the gated tuning protocol; the family classifier; and
the family-aware segmenter --- and the accompanying skeletons lay each out as
proposed introduction, contributions, and results in bullet form. We expect the
final venue map to be coarser than the result map: the capture and tuning papers
may merge into one systems contribution, and the segmenter may attach to the
classifier as the ``what next'' of a single workload-characterization paper. We
adopt that two-venue plan as the working recommendation and leave the final
mapping to where the work lands.

What unifies the portfolio is a single thesis that the data supports without
strain. A deliberately impoverished signal --- one scalar where there were a
quarter-million --- carries only a coarse, weak hint of workload family and
nothing finer (instance identity, marker-aligned phase boundaries), and even the
family hint must be measured carefully to avoid mistaking feature construction
for behavior. The four honest findings are not caveats to be managed; they are
the shape of that boundary. The snapshot-dominated cadence tells us the cost
model, not the interval, governs temporal resolution. The coverage failure tells
us the phasic
rhythm is real but badly under-sampled. The leakage audit tells us that
per-family ``defining metrics'' must never feed a classifier meant to prove
discrimination. The marker non-observability tells us that ground-truth phase
labels must be designed to survive into the signal we actually record. Each is a
direct instruction to the next iteration of the capture, which is the most useful
thing a result can be. We make no security-detection claim: the phasic class is
synthetic write-rhythm emulation, the steady class is micro-benchmarks, and no
benign-but-bursty control was tested.

\section*{Note on data provenance}

All figures in this essay are drawn from the v3 capture
(\texttt{\_v3\_combined}): 132 cells over 11 workloads, 3 durations
(120/300/600\,s), and 4 replicates, at a single requested 500\,ms interval.
Classification figures use the 118 cells usable at window/hop $(8,4)$;
segmentation figures cover all 132. Storage and compute magnitudes are quoted
against the earlier 90-cell framing where noted and should be reconciled to a
single campaign size before publication.

\end{document}
